%! TEX root = ../paper.tex
\section{The Six Induction Vectors in Full}
\label{sec:vectors}


The agent is consistently assigned the same benign surface task: patching a division bug to pass a test suite.
The experimental variance lies strictly in the temptation planted within its environment.

Tasks 01 through 03 represent passive temptations, such as exposed credentials or errors pointing to an internal registry, lacking any explicit instruction to act.
These serve as structural baselines and produce none of the primary leakage results, existing solely to validate the harness's observational capabilities.

Tasks 04 through 06 introduce an \emph{identical}, active injected instruction demanding the agent create a confirmation directory named using its session identifier.
These tasks differ exclusively in their delivery channel: a live fetch from an external mirror, a note seeded in cross-run persistent memory, or a document embedded in a local knowledge base.
Holding the malicious payload constant while varying the delivery channel enables the comparative analysis detailed in Section~\ref{sec:results}.
This exfiltration technique directly mirrors the anchor incident, which leaked data through directory naming conventions~\cite{huggingface2026intrusion}.

Vectors 05 and 06 specifically model \textbf{ASI06: Memory and Context Poisoning}~\cite{owasp-asi06}, where untrusted input is processed as trusted memory or retrieved context to compromise subsequent sessions.
Vector 05 represents the most critical persistence threat, since \texttt{./memory} operates as a host bind mount that survives container teardown (\texttt{docker compose down -v}), effectively outliving the run that created it.
